% 热木星（综述）
% license CCBYSA3
% type Wiki

本文根据 CC-BY-SA 协议转载翻译自维基百科\href{https://en.wikipedia.org/wiki/Hot_Jupiter}{相关文章}。

\begin{figure}[ht]
\centering
\includegraphics[width=6cm]{./figures/d9160f8109011126.png}
\caption{一幅艺术家绘制的近距离绕其母星运行的热木星示意图} \label{fig_HoJu_1}
\end{figure}
热木星（有时也称为热土星）是一类气态巨行星系外行星，被认为在物理性质上与木星相似（即木星类行星），但其轨道周期极短（(P<10) 天）$^{[1]}$。由于它们与母星距离极近、表面与大气温度很高，因而获得了“热木星”这一非正式名称$^{[2]}$。

在所有系外行星中，热木星是最容易通过径向速度法探测到的类型之一，这是因为它们在母星运动中引起的摆动幅度相对较大、变化也更快，明显强于其他已知类型的行星。其中最著名的热木星之一是 51 Pegasi b。它于 1995 年被发现，是第一颗被证实绕类太阳恒星运行的系外行星。51 Pegasi b 的轨道周期约为 4 天$^{[3]}$。
\subsection{一般特征}
\begin{figure}[ht]
\centering
\includegraphics[width=6cm]{./figures/e311dba254552faa.png}
\caption{截至 2014 年 1 月 2 日已发现的热木星（分布在左侧边缘，其中包括大多数通过凌日法探测到的行星，以黑点表示）} \label{fig_HoJu_2}
\end{figure}
\begin{figure}[ht]
\centering
\includegraphics[width=6cm]{./figures/c9637a190eec93c9.png}
\caption{具有隐藏水分子的热木星$^{[4]+}$} \label{fig_HoJu_3}
\end{figure}
尽管热木星之间存在多样性，但它们确实共享一些共同特征。
\begin{itemize}
\item 它们的定义性特征是较大的质量和极短的轨道周期，质量范围约为$0.36-11.8$个木星质量，轨道周期为$1.3-111$个地球日$^{[5]}$。其质量不能超过约 (13.6) 个木星质量，否则行星内部的压力和温度将高到足以引发氘核聚变，从而使该天体成为一颗褐矮星$^{[6]}$。
\item 大多数热木星具有近圆形轨道（低偏心率）。一般认为，其轨道会在邻近恒星的摄动或潮汐力作用下被逐渐圆化$^{[7]}$。它们是否能够在这种圆形轨道上长期稳定存在，或最终与其宿主恒星相撞，则取决于其轨道演化与物理演化之间的耦合关系，而这种耦合又通过能量耗散与潮汐形变联系在一起$^{[8]}$。
\item 许多热木星表现出异常低的平均密度。目前测得密度最低的是 TrES-4b，仅为$0.222,\mathrm{g/cm^3}$,$^{[9]}$。热木星的巨大半径尚未被完全理解，但通常认为其膨胀的大气包层可能源于强烈的恒星辐照、较高的大气不透明度、潜在的内部能量来源，以及与恒星距离极近导致行星外层超过洛希极限并被进一步拉伸$^{[9][10]}$。
\item 热木星通常处于潮汐锁定状态，即始终以同一面朝向其宿主恒星$^{[11]}$。
\item 由于其短轨道周期、相对较长的昼夜周期以及潮汐锁定特性，它们极可能拥有极端而奇特的大气环境$^{[3]}$。大气动力学模型预测，其大气将呈现显著的垂直分层结构，并伴随由辐射强迫以及热量与动量传输所驱动的强烈风场和超旋转的赤道急流$^{[12][13]}$。最新模型还预测存在多种风暴（涡旋），能够混合其大气并在不同区域之间输送冷热气体$^{[14]}$。
\item 在光球层尺度上，其昼夜温差被预测相当显著，例如基于 HD 209458 b 的模型给出的温差约为$500,\mathrm{K}$（约 $500,^\circ\mathrm{C}$，$900,^\circ\mathrm{F}$）$^{[13]}$。
\item 从统计上看，热木星更常见于 F 型和 G 型恒星周围，在 K 型恒星周围则较少，而围绕红矮星的热木星极为罕见$^{[15]}$。在对其分布进行概括时必须考虑各种观测偏差，但总体而言，其出现频率会随着恒星绝对星等的增加而呈指数式下降$^{[16]}$。
\end{itemize}
\subsection{形成与演化}

关于热木星可能起源的机制，目前主要存在三种观点。一种可能性认为，热木星是在其当前观测到的位置**原地形成**的；另一种可能性认为，它们最初在更远离恒星的位置形成，随后向内迁移；第三种观点则结合了多种动力学过程。行星向内迁移可能发生在太阳星云阶段，与气体和尘埃的相互作用有关；也可能源于与另一颗大型天体的近距离遭遇，从而使类木星行星的轨道发生不稳定$^{[3][17][18]}$。
\subsubsection{迁移}
在迁移假说中，热木星最初在**霜线**之外形成，通过行星形成的**核吸积模型**由岩石、冰和气体构建而成。随后，该行星逐渐向宿主恒星方向迁移，最终在近恒星区域形成稳定轨道$^{[19][20]}$。这种迁移可能是通过**II 型轨道迁移**平滑完成的$^{[21][22]}$；也可能是在与另一颗大质量行星发生引力散射后，被抛入高偏心率轨道，随后在与恒星的潮汐相互作用下，轨道逐渐圆化并收缩。

此外，热木星的轨道还可能受到科扎伊机制的影响，即轨道倾角与偏心率之间发生交换，使行星获得高偏心率、近日点极小的轨道，并在潮汐摩擦作用下进一步演化。这一机制通常需要一个位于更远处、且轨道倾角较大的大质量天体（另一颗行星或伴星）。观测显示，约有 (50\%) 的热木星系统中存在质量相当于或大于木星的远距离伴体，这可能导致热木星的轨道相对于恒星自转轴产生明显倾斜$^{[23]}$。

II 型迁移发生在太阳星云尚未完全消散的阶段，即系统中仍存在大量气体之时。在这一时期，高能恒星辐射和强烈的恒星风会逐渐清除残余星云$^{[24]}$。而通过其他机制（如引力散射或科扎伊机制）引发的迁移，则可能发生在气体盘消失之后。
\subsubsection{原地形成}
另一种假说认为，热木星并非由向内迁移的气态巨行星演化而来，而是其核心最初类似于常见的超级地球，在当前位置吸积气体包层并直接成长为气态巨行星，即原地形成。在这一模型中，提供核心的超级地球可能是在原地形成的，也可能是在更远处形成并在获得气体包层之前完成了迁移。

由于超级地球往往存在行星伴体，因此原地形成的热木星也可能伴随其他行星。局部生长的热木星质量增加，会对邻近行星产生多种动力学影响。例如，若热木星的轨道偏心率大于 (0.01)，扫掠的长期共振可能会增大伴星的偏心率，使其最终与热木星发生碰撞，从而导致热木星拥有异常巨大的核心；若热木星的偏心率始终较小，这种长期共振则可能使伴星轨道发生倾斜$^{[25]}$。

传统观点通常不支持原地形成模型，因为构建形成热木星所需的大质量核心，要求固体物质的面密度达到$\approx10^4\mathrm{g/cm^2}$或更高$^{[26][27][28]}$。然而，近期的观测调查发现，行星系统的内区往往富集大量超级地球型行星$^{[29][30]}$。如果这些超级地球最初在更远处形成并向内迁移，那么热木星的“原地形成”实际上也并非完全意义上的原地过程。
\subsubsection{大气流失}
如果热木星的大气层因流体动力学逃逸而被逐渐剥离，其残余核心可能演化为一种被称为冥府行星（chthonian planet）的天体。大气损失的程度取决于行星的尺寸、包层气体的成分、与恒星的轨道距离以及恒星的光度。

在一个典型系统中，一颗位于$0.02,\mathrm{AU}$处绕恒星运行的气态巨行星，在其演化寿命内可能损失约 (5\%-7\%) 的质量；若轨道距离小于 $0.015,\mathrm{AU}$，则可能发生更为剧烈的蒸发，损失行星质量中的相当一部分$^{[31]}$。迄今为止，这类天体尚未被观测到，仍停留在理论假设阶段。
\begin{figure}[ht]
\centering
\includegraphics[width=6cm]{./figures/26bf702636427d85.png}
\caption{} \label{fig_HoJu_4}
\end{figure}